

\begin{derivation}[从\eqnrefs{eq:sm-covariant-derivative-dp11}到\eqnrefs{eq:sm-field-strengths-dp11}]
从正文\eqnrefs{eq:sm-covariant-derivative-dp11} 的协变条件出发，先求规范联络的有限变换律，再计算两个协变导数的对易子，并分别代入三个规范群因子。

把三个群因子分别写成
\begin{equation*}
 U_3(x)=\exp[\ii\alpha^a(x)T^a],\qquad
 U_2(x)=\exp[\ii\theta^I(x)t^I],\qquad
 U_Y(x)=\exp[\ii Y\vartheta_Y(x)].
\end{equation*}
对表示 $(R_3,R_2,Y)$ 中的场，局域变换为
\begin{equation*}
 \Psi'(x)=U_3(x)U_2(x)U_Y(x)\Psi(x).
\end{equation*}
三个群因子作用在颜色、弱同位旋和超荷三个不同内部空间，因而彼此交换。规范联络的变换律可先在一个一般非阿贝尔因子上推导，再分别代入 $SU(3)_c$ 与 $SU(2)_L$；阿贝尔 $U(1)_Y$ 是同一结果在对易生成元下的特例。记
\begin{equation*}
 D_\mu=\partial_\mu+\ii g_{\rm G}\mathscr C_\mu,
 \qquad \mathscr C_\mu=\mathscr C_\mu^aT^a,
\end{equation*}
并要求协变导数本身按相应表示协变变换：
\begin{equation*}
 D_\mu'\Psi'=U D_\mu\Psi.
\end{equation*}
左边逐项展开，
\begin{align*}
 D_\mu'\Psi'
 &=(\partial_\mu+\ii g_{\rm G}\mathscr C_\mu')U\Psi,\\
 &=(\partial_\mu U)\Psi+U\partial_\mu\Psi
   +\ii g_{\rm G}\mathscr C_\mu'U\Psi.
\end{align*}
右边为
\begin{equation*}
 U D_\mu\Psi
 =U\partial_\mu\Psi+\ii g_{\rm G}U\mathscr C_\mu\Psi.
\end{equation*}
消去共同的 $U\partial_\mu\Psi$ 后，任意 $\Psi$ 均须满足
\begin{equation*}
 (\partial_\mu U)+\ii g_{\rm G}\mathscr C_\mu'U
 =\ii g_{\rm G}U\mathscr C_\mu.
\end{equation*}
右乘 $U^{-1}$，得到连接的有限局域变换律
\begin{equation*}
 {
 \mathscr C_\mu'
 =U\mathscr C_\mu U^{-1}
 +\frac{\ii}{g_{\rm G}}(\partial_\mu U)U^{-1}.}
\end{equation*}
第二项是非齐次项；它正好抵消普通导数作用在局域群元上产生的 $\partial_\mu U$。若 $U$ 与 $\mathscr C_\mu$ 交换，上式退化为阿贝尔规范联络的平移。

对超荷群，直接取 $U_Y=\exp(\ii Y\vartheta_Y)$。由
\begin{align*}
 D_\mu'\Psi'
 &=\left(\partial_\mu+\ii g'Y\mathcal B_\mu'\right)
   \ee^{\ii Y\vartheta_Y}\Psi,\\
 &=\ee^{\ii Y\vartheta_Y}
 \left[\partial_\mu+\ii Y\partial_\mu\vartheta_Y
 +\ii g'Y\mathcal B_\mu'\right]\Psi
\end{align*}
与 $D_\mu'\Psi'=\ee^{\ii Y\vartheta_Y}D_\mu\Psi$ 比较，得到
\begin{equation*}
 {
 \mathcal B_\mu'=\mathcal B_\mu-\frac1{g'}\partial_\mu\vartheta_Y.}
\end{equation*}

接着由协变导数的对易子推出曲率，也就是规范场强。对任意检验场 $\Psi$，
\begin{align*}
 [D_\mu,D_\nu]\Psi
 =&\left[\partial_\mu+\ii g_{\rm G}\mathscr C_\mu,
 \partial_\nu+\ii g_{\rm G}\mathscr C_\nu\right]\Psi,\\
 =&\ii g_{\rm G}(\partial_\mu \mathscr C_\nu-\partial_\nu \mathscr C_\mu)\Psi
 -(g_{\rm G})^2[\mathscr C_\mu,\mathscr C_\nu]\Psi.
\end{align*}
利用
\begin{equation*}
 [\mathscr C_\mu,\mathscr C_\nu]
 =\mathscr C_\mu^a\mathscr C_\nu^b[T^a,T^b]
 =\ii f^{abc}\mathscr C_\mu^a\mathscr C_\nu^bT^c,
\end{equation*}
可整理为
\begin{equation*}
 [D_\mu,D_\nu]
 =\ii g_{\rm G}\mathcal F_{\mu\nu},
\end{equation*}
其中
\begin{equation*}
 {
 \mathcal F_{\mu\nu}^a
 =\partial_\mu \mathscr C_\nu^a-\partial_\nu \mathscr C_\mu^a
 -g_{\rm G}f^{abc}\mathscr C_\mu^b\mathscr C_\nu^c.}
\end{equation*}
由于上式等价于算符恒等式 $D_\mu'=UD_\mu U^{-1}$，
\begin{align*}
 [D_\mu',D_\nu']
 &=UD_\mu U^{-1}UD_\nu U^{-1}
  -UD_\nu U^{-1}UD_\mu U^{-1},\\
 &=U[D_\mu,D_\nu]U^{-1},
\end{align*}
从而
\begin{equation*}
 {\mathcal F_{\mu\nu}'=U\mathcal F_{\mu\nu}U^{-1}.}
\end{equation*}
规范联络含非齐次项，而场强按伴随表示齐次变换，因此 $\Tr(\mathcal F_{\mu\nu}\mathcal F^{\mu\nu})$ 具有规范不变性。把这一一般结果分别取 $(g_{\rm G},\mathscr C_\mu,T^a,f^{abc})=(g_s,\mathscr C_\mu,T^a,f^{abc})$ 与 $(g,\mathcal W_\mu,t^I,\epsilon^{IJK})$，便得到色规范场与弱规范场的场强；阿贝尔超荷群满足 $f^{abc}=0$，所以场强只保留普通旋度项。
\end{derivation}
